% Generated draft: requires mathematical and source review.
\documentclass[11pt]{article}
\usepackage[margin=25mm]{geometry}
\usepackage{fontspec}
\setmainfont{Latin Modern Roman}
\usepackage{amsmath,amssymb,mathtools,unicode-math}
\setmathfont{Latin Modern Math}
\usepackage{enumitem,needspace,xurl,hyperref}
\setlength{\parindent}{0pt}
\setlength{\parskip}{5pt}
\setlength{\emergencystretch}{3em}
\newcommand{\N}{\mathbb{N}}
\newcommand{\Z}{\mathbb{Z}}
\newcommand{\Q}{\mathbb{Q}}
\newcommand{\R}{\mathbb{R}}
\newcommand{\C}{\mathbb{C}}
\newcommand{\F}{\mathbb{F}}
\newcommand{\A}{\mathbb{A}}
\newcommand{\PP}{\mathbb{P}}
\newcommand{\E}{\mathbb{E}}
\newcommand{\eps}{\varepsilon}
\newcommand{\dd}{\,\mathrm{d}}
\begin{document}
\section*{UnsolvedMath ID: 30001232}
\textbf{Source problem code:} catalogue entry
\subsection{Definitions and mathematical statement}
For the usual smooth-surface interpretation, let \(S\) be a smooth projective surface over an algebraically closed field, let \(L\) be an ample line bundle on \(S\), and let \(x\) be a closed point of \(S\). Define the Seshadri constant by
\[
\varepsilon(S,L;x)=\inf_{\substack{C\subseteq S\text{ reduced irreducible curve}\\x\in C}}\frac{L\cdot C}{\operatorname{mult}_x(C)},
\]
where \(L\cdot C\) is the intersection number and \(\operatorname{mult}_x(C)\) is the multiplicity of \(C\) at \(x\). Let \(K_S\) denote a canonical divisor and let \(K_S^2=K_S\cdot K_S\). In this interpretation, minimality means that \(S\) contains no smooth rational curve of self-intersection \(-1\).

The supplied question asks whether, for every minimal projective surface \(S\), every ample line bundle \(L\) on \(S\), and every point \(x\in S\), one has
\[
\varepsilon(S,L;x)\geq\frac{1}{\sqrt[4]{\,2+|K_S^2|\,}}.
\]
The fourth root is the positive real fourth root. The bound includes equality, is required at every point, and has no restriction on the Picard number.

The source record does not specify the ground field, its characteristic, or smoothness. Thus the definitions above give a conditional smooth-surface interpretation; they do not establish that this is the full intended scope. In particular, an interpretation allowing singular surfaces would require additional conventions for minimality and \(K_S^2\).
\subsection{Short English statement}
The question asks whether every ample line bundle on a minimal projective surface has, at every point, a Seshadri constant at least \((2+|K_S^2|)^{-1/4}\), with arbitrary Picard number. The ground field and smoothness assumptions require clarification. The supplied open status is a dated curation label.
\subsection{Sources}
\begin{itemize}
\item Local registry: supplied UnsolvedMath record, permanent ID \(30001232\), title \emph{Lower Bounds for Seshadri Constants on Minimal Surfaces}, exact source problem code \texttt{catalogue entry}. The record belongs to the Oberwolfach Reports Open Problems collection and gives the proposed year as 2009; no individual proposer is supplied.
\item Supplied report citation and literature source labelled ``original report'': \emph{Multiplier Ideal Sheaves in Algebraic and Complex Geometry} (2009), Oberwolfach Reports, pp.\ 1101--1156, DOI 10.4171/owr/2009/21. \url{https://doi.org/10.4171/owr/2009/21}. The supplied background describes Oberwolfach Reports as documenting mathematical workshops and their open-problem discussions.
\item Supplied literature source labelled ``current Seshadri survey'': \url{https://www.sciencedirect.com/science/article/pii/S0007449725000430}. No further bibliographic details are supplied.
\item Literature assessment dated exactly \texttt{2026-08-21}: the supplied curation labels the problem open and reports that no general proof or disproof was found for the displayed lower bound on minimal surfaces with arbitrary Picard number. This is supplied dated curation derived from the OWR record, not an independently verified literature assessment; the time-sensitive status and cited claims require independent verification.
\item The registry supplies statement status \texttt{corrected\_verified}, with the explanation that the proposed lower bound is complete and that surface and root notation were normalized. This is a supplied verification label, not independent mathematical verification here.
\end{itemize}
\end{document}
